\documentclass[10pt,a4paper]{amsart}
\usepackage[margin=19mm]{geometry}
\usepackage{amsmath,amssymb,mathtools}
\usepackage[scr=rsfs,cal=euler]{mathalfa}
\usepackage{xfrac}
\usepackage[unicode,hidelinks,bookmarks=false]{hyperref}
\newcommand{\ie}{i.e.\ }
\newcommand{\cf}{cf.\ }
\newcommand{\resp}{respectively\ }
\newcommand{\Subsection}{\subsection}
\providecommand{\nobreakdash}{\nobreak\hbox{-}}
\setlength{\emergencystretch}{2em}
\multlinegap0pt
\usepackage{amssymb}
\usepackage{bm}
\usepackage[shortlabels]{enumitem}
\newtheorem{proposition}{Proposition}
\newtheorem{lemma}{Lemma}
\newcommand{\Id}{\operatorname{Id}}
\newcommand{\Sph}{\mathbb S}
\DeclareMathOperator{\tr}{tr}
\title{Free boundary flows by powers of the Gauss curvature in the unit ball}
\author{Tianci Luo, Yong Wei, Rong Zhou}
\date{Source: arXiv:2607.26923v1 (CC BY 4.0)}
\begin{document}
\maketitle

\noindent\emph{Source and licence.} Free boundary flows by powers of the Gauss curvature in the unit ball. Version arXiv:2607.26923v1 (CC BY 4.0), distributed by arXiv under the Creative Commons Attribution 4.0 International licence, http://creativecommons.org/licenses/by/4.0/, as recorded in the arXiv licence metadata for this exact version and shown on the abstract page https://arxiv.org/abs/2607.26923v1. Full licence notice: \texttt{licenses/arxiv-2607.26923-CC-BY-4.0.txt}.\par\smallskip

\noindent\textit{Editorial scope.} Counted unit: the article's lemma prop:F-lower-alpha (source0.tex lines 2636--2642). The article's complete original proof of it is delivered with it (source0.tex lines 2644--2883, one \texttt{proof} environment of 240 lines). No mathematics is written, repaired or re-derived: the statement and the proof are the article's own lines, in the article's own notation, and no retained line is substituted or re-flowed.  Retained uncounted as the source-local context the proof needs: lines 544--547; lines 2377--2389; lines 2494--2496; lines 2559--2565.  Nothing was omitted from any retained span, so the package carries no omission record.  No retained line contains a citation, so the wrapper carries no bibliography and no bibliography entry.  No retained line contains a figure environment, an image-inclusion command or drawing code, so no image is delivered and the package declares no asset dependency.  Editorial formatting only, in addition: an AMS \texttt{amsart} preamble that restates the article's own declarations and macros used by the retained lines, namely the environments \texttt{proposition}, \texttt{lemma} on separate counters, together with the standard packages those lines need.  The retained environments print in this document as Proposition 1, Lemma 1 in order of appearance, whereas the article prints the same items with the numbering of its own full document.  The complete native source of the article as distributed in the arXiv e-print for this exact version is bundled unchanged as \texttt{sources/206304/source0.tex}.
\begin{equation}\label{eq:kappa-lower-alpha}
        \kappa_i(x,t)\ge c_1,
        \qquad i=1,\ldots,n,
\end{equation}

\begin{proposition}
	\label{thm:C0-alpha}
	Assume $\alpha>\frac{1}{n+2}$. Then
	\begin{equation}
    \label{eq:entropy-point-to-zero}
		z_e(s)\longrightarrow 0.
	\end{equation}
	Consequently, after increasing the initial time, the support function satisfies
	\begin{equation}\label{eq:subunit-origin-C0}
		c\le h(\xi,s)\le C,\qquad |\overline{\nabla}h|_{\sigma}(\xi,s)\leq C,
		\qquad \xi\in\mathbb{S}^n_+.
	\end{equation}
\end{proposition}

\begin{equation}\label{eq:F-shifted}
    \partial_s h=h-\frac{W}{\zeta},
\end{equation}

\begin{lemma}\label{lem:Lambda-bounds-alpha}
There is $C\ge1$ such that
\begin{equation}\label{eq:Lambda-bounds-alpha}
        C^{-1}\le \zeta(s)= \dfrac{1}{|\mathbb{S}_+^n|}\int_{\mathbb{S}_+^n} \Theta K^{\alpha-1}\mathrm{d}\sigma  \le C
\end{equation}
for all sufficiently large $s$.
\end{lemma}

\begin{lemma}[Lower speed bound]\label{prop:F-lower-alpha}
There is $c>0$ such that
\begin{equation}\label{eq:F-lower-alpha}
       W(\xi,s)\ge c
\end{equation}
for all $s$ sufficiently large.
\end{lemma}

\begin{proof}
Recall that
\begin{equation*}
    \omega=\rho q,
    \qquad
    \widehat B=\widehat A^{-1},
    \qquad
    P=A\widehat B A.
\end{equation*}
We define the operator
\begin{equation*}
        \mathcal L
        =\partial_s-\frac{\alpha W}{\zeta}P^{ij}\overline{\nabla}_i\overline{\nabla}_j.
\end{equation*}

We first compute $\mathcal{L}W$. Using \eqref{eq:F-shifted} we obtain
\begin{equation*}
        \partial_sB_{ij}
        =\overline{\nabla}_i\overline{\nabla}_j \partial_sh+\partial_sh\sigma_{ij}
        =B_{ij}-\frac{1}{\zeta}(W_{ij}+W\sigma_{ij}).
\end{equation*}
Consequently
\begin{equation*}
        \partial_sA=-A(\partial_sB)A
        =-A+\frac{1}{\zeta} A(\overline{\nabla}^2 W+ W\sigma)A,
\end{equation*}
and therefore
\begin{equation}\label{eq:hatA-evol-lower-alpha}
        \partial_s\widehat A
        =-A+\frac{1}{\zeta} A(\overline{\nabla}^2 W+ W\sigma)A+(\partial_s\omega)\Id .
\end{equation}
Differentiating
\begin{equation*}
        \log W=(n\alpha+1)\log \Lambda+\alpha\log\widehat K
\end{equation*}
and using \eqref{eq:hatA-evol-lower-alpha}, we obtain the exact identity
\begin{align}
        \frac{\mathcal L W}{W}
        &=
        (n\alpha+1)\,\partial_s\log \Lambda
        -\alpha\tr(\widehat B A)
        +\frac{\alpha W}{\zeta}\tr P
        +\alpha(\partial_s\omega)\tr\widehat B   \nonumber \\
        &=
        -n\alpha
        +(n\alpha+1)\,\partial_s\log \Lambda
        +\alpha(\partial_s\omega+\omega)\tr\widehat B
        +\frac{\alpha W}{\zeta}\tr P,   \label{eq:LF-rough}
\end{align}
where in the last equality we used
\begin{equation*}
        A=\widehat A-\omega\Id,\qquad
        \tr(\widehat B A)=n-\omega\tr\widehat B .
\end{equation*}
Similarly, using $h_{ij}=B_{ij}-h\sigma_{ij}$ we obtain
\begin{align}
        \mathcal{L} h
        &=h-\frac{W}{\zeta}
          -\frac{\alpha W}{\zeta}P^{ij}h_{ij} \nonumber  \\
        &=h-\frac{W}{\zeta}
          -\frac{\alpha W}{\zeta}P^{ij}B_{ij}
          +\frac{\alpha W}{\zeta}h\tr P     \nonumber     \\
        &=h-\frac{W}{\zeta}
          -\frac{\alpha W}{\zeta}\tr(\widehat B A)
          +\frac{\alpha W}{\zeta}h\tr P.  \label{eq:Lh-exact-lower-alpha}
\end{align}

Fix $\ell>0$ and consider the function $\log(Wh^\ell)$. At its spatial minimum,  we have
\begin{equation}\label{eq:min-condition-Fh-alpha}
        \frac{W_i}{W}+\ell\frac{h_i}{h}=0 .
\end{equation}
Using \eqref{eq:LF-rough} and \eqref{eq:Lh-exact-lower-alpha} we obtain at such a point
\begin{align}
       \mathcal L\log(Wh^\ell)
        &=
        \frac{\mathcal L W}{W}
        +\ell\frac{\mathcal L h}{h}
        +\frac{\alpha W}{\zeta}P^{ij}
          \left(
             \frac{W_iW_j}{W^2}
             +\ell\frac{h_i h_j}{h^2}
          \right)                        \nonumber              \\
        &=
        \ell-n\alpha
        +(n\alpha+1)\,\partial_s\log \Lambda
        +\alpha(\partial_s\omega+\omega)\tr\widehat B    \nonumber   \\
        &\quad
        -\ell\frac{W}{\zeta h}
        -\ell\frac{\alpha W}{\zeta h}\tr(\widehat B A)
        +(\ell+1)\frac{\alpha W}{\zeta}\tr P        \nonumber      \\
        &\quad
        +\frac{\alpha W}{\zeta}P^{ij}
          \left(
             \frac{W_iW_j}{W^2}
             +\ell\frac{h_i h_j}{h^2}
          \right).           \label{eq:LlogFh-exact-alpha}
\end{align}
The last line is nonnegative.  The term involving $\tr P$ is also
nonnegative and will be discarded from below.

It remains to estimate the terms coming from $\Lambda$ and $\omega=\rho q$.
At the minimum point of $\log(Wh^\ell)$, \eqref{eq:min-condition-Fh-alpha} gives
\begin{equation*}
        \overline{\nabla} \partial_s h
        =\overline{\nabla} h-\frac{\overline{\nabla} W}{\zeta}
        =\left(1+\ell\frac{W}{\zeta h}\right)\overline{\nabla} h.
\end{equation*}
Since $h$ and $\overline{\nabla} h$ are uniformly bounded by the $C^0$ estimate and convexity, we have
\begin{equation}
    |\partial_s h|+|\overline{\nabla}\partial_s h|_{\sigma}\leq C\left( 1+\dfrac{W}{\zeta h} \right).
    \label{s6.partialshbd}
\end{equation}
Recalling the definitions
\begin{align*}
        \Lambda=&1+2\rho \langle \overline{\nabla}h+h\xi,E \rangle+\rho^2(h^2+|\overline{\nabla}h|_{\sigma}^2), \\
        q=&-\frac{2(\langle \xi,E \rangle+\rho h)}{\Lambda},\qquad \rho_s=-\rho,
\end{align*}
and using the $C^0$ and $C^1$ estimates, \eqref{s6.partialshbd}, and $\Lambda=1+O(\delta(s))$ at the same point, we get
\begin{align}
    |\partial_s\log \Lambda| %=& -2\rho\langle \overline{\nabla}h+h\xi,E \rangle + 2\rho \langle \overline{\nabla}\partial_s h+\partial_s h\xi,E \rangle \nonumber\\
    %&-2\rho^2 (|\overline{\nabla}h|_{\sigma}^2+h^2)+2\rho^2\left(h\partial_s h +  \overline{\nabla}_i h\overline{\nabla}_i(\partial_s h)\right) \nonumber\\
    \le C\rho\left(1+\frac{W}{\zeta h}\right).  \label{eq-Lambdabd}
\end{align}
Therefore, combining the $C^0$ and $C^1$ estimates with $\Lambda=1+O(\delta(s))$, \eqref{s6.partialshbd}, and \eqref{eq-Lambdabd}, yields
\begin{align*}
        |\,\partial_s\omega+\omega\,|
        =&\rho|\partial_s q|\\
        =&  \rho\left| \dfrac{2\rho}{\Lambda}(h-\partial_s h)-q\partial_s\log\Lambda \right|\\
        \le & C \rho^2 \left(1+\frac{W}{\zeta h}\right).
\end{align*}
The conformal relation $\widehat A=\rho e^w\overline A$ and the lower curvature bound \eqref{eq:kappa-lower-alpha} for $\overline{A}$ give
\begin{equation*}
        \widehat A\ge c\rho\Id , 
        \qquad
        \rho\tr\widehat B\le C.
\end{equation*}
Also,
\begin{equation}
        \tr(\widehat B A)
        =n-\omega\tr\widehat B,
        \qquad
        |\omega|\tr\widehat B\le C.
        \label{s6.BA}
\end{equation}
Combining \eqref{eq-Lambdabd}-\eqref{s6.BA} with \eqref{eq:LlogFh-exact-alpha} gives, at a spatial minimum of $\log(Wh^\ell)$,
\begin{equation}\label{eq:logFhl}
        \mathcal L\log(Wh^\ell)
        \ge
        \ell-n\alpha
        -C_1\rho
        -C_\ell\frac{W}{\zeta h}.
\end{equation}


Choose $\ell=n\alpha+2$ and $s_3$ sufficiently large that $C_1\rho\le 1$.  By Proposition \ref{thm:C0-alpha} and Lemma \ref{lem:Lambda-bounds-alpha}, there is a
constant  $\Lambda_0>1$ such that
\begin{equation*}
        \zeta\ge\Lambda_0^{-1},
        \qquad
        \Lambda_0^{-1}\le h\le\Lambda_0
\end{equation*}
for all $s\ge s_3$. Then \eqref{eq:logFhl} implies
\begin{equation*}
    \mathcal L\log(Wh^{\ell})\ge 1-C_2 W
\end{equation*}
at a spatial minimum point.



Choose $\lambda>0$ so small that
\begin{equation*}
        \lambda < \frac{1}{4C_2\Lambda_0^{\ell}}
\end{equation*}
and, after decreasing $\lambda$ if necessary, $ Wh^\ell(\cdot,s_3)\ge 2\lambda$.  

We claim that
\begin{equation*}
        Wh^\ell\ge\lambda
        \qquad\hbox{on }\Sph^n_+\times[s_3,\infty).
\end{equation*}
Suppose otherwise.  Let $s_*>s_3$ be the first time at which the minimum of
$Wh^\ell$ equals $\lambda$, and let $\xi_*$ be a minimum point.  
We observe that the following maximum principle argument remains
valid if $\xi_*$ lies on the boundary. Indeed, differentiating the
boundary condition $\overline{\nabla}_{\eta}h=0$ with respect to $s$
and using \eqref{eq:F-shifted}  
gives $\overline{\nabla}_{\eta}W=0$ on $\partial\mathbb S^n_+$.
Consequently,
\begin{equation*}
    \overline{\nabla}_{\eta}\log(Wh^\ell)=0
    \qquad\text{on }\partial\mathbb S^n_+.
\end{equation*}
Since the equator is totally geodesic and
$\overline{\nabla}_{\eta}h=0$, we also have
\begin{equation*}
    B_{\alpha\eta}=0
    \qquad\text{on }\partial\mathbb S^n_+.
\end{equation*}
It follows that $A$, $\widehat A$, $\widehat B$, and
$P=A\widehat BA$ are block diagonal with respect to the decomposition
into the tangential and conormal directions. In particular, $P^{\alpha\eta}=0$. 
At a boundary minimum of $\log(Wh^\ell)$, the one-sided second derivative
test gives
\begin{equation*}
    \left(\overline{\nabla}_{\alpha}
    \overline{\nabla}_{\beta}\log(Wh^\ell)\right)\geq0,
    \qquad
    \overline{\nabla}_{\eta}
    \overline{\nabla}_{\eta}\log(Wh^\ell)\geq0.
\end{equation*}
Therefore,
\begin{equation*}
    P^{ij}\overline{\nabla}_i
    \overline{\nabla}_j\log(Wh^\ell)\geq0.
\end{equation*}
Since $s_*$ is the first time at which the minimum of $Wh^\ell$
reaches $\lambda$, we also have
\begin{equation*}
    \partial_s\log(Wh^\ell)(\xi_*,s_*)\leq0.
\end{equation*}
Thus, whether $\xi_*$ is an interior point or a boundary point,
\begin{equation*}
    \mathcal L\log(Wh^\ell)(\xi_*,s_*)\leq0.
\end{equation*}
On the other hand, at $(\xi_*,s_*)$,
\begin{equation*}
    0
    \geq \mathcal L\log(Wh^\ell)
    \geq 1-C_2W
    =1-C_2\frac{\lambda}{h^\ell}
    \geq 1-\frac14
    >0,
\end{equation*}
which is a contradiction. Hence $Wh^\ell\ge\lambda$.  Since $h\le\Lambda_0$,
we obtain
\begin{equation*}
        W\ge \lambda\Lambda_0^{-\ell}=:c>0.
\end{equation*}
This proves \eqref{eq:F-lower-alpha}.
\end{proof}

\end{document}
